%! Characteristics of Polynomial Functions — Unit 3, Lesson 3.0 — Group Activity (Answer Key)
\documentclass[10pt]{article}
\usepackage{algebra2-article}
\usepackage{algebra2-key}   % pulls in -boxes; do NOT also load -boxes
\newcommand{\TallMath}[1]{$\displaystyle #1\rule[-1.4em]{0pt}{3.2em}$}
\newcommand{\lgrid}[4]{%
  \draw[step=1, linegray!40, very thin] (#1,#3) grid (#2,#4);
  \draw[->, semithick, charcoal] (#1,0)--(#2,0) node[below right, font=\scriptsize]{$x$};
  \draw[->, semithick, charcoal] (0,#3)--(0,#4) node[left, font=\scriptsize]{$y$};}

\begin{document}

\pageheader{Unit 3, Lesson 3.0}{Group Activity --- Answer Key}
\namepartnerperiod

\vspace{0.08in}

\begin{headlinebox}{forestbg}
\small\textbf{Reading Polynomial Graphs.} Every characteristic of a polynomial is something you can
\emph{point to} on its graph. With your partner, read each graph carefully and \emph{justify} your
answers --- especially the end behavior, the multiplicity (cross vs.\ touch), and relative vs.\
absolute extrema.
\end{headlinebox}

\vspace{0.06in}

\begin{tcolorbox}[colback=white, colframe=black!40, title=\textbf{Tier R --- Remediate},
    fonttitle=\bfseries, arc=1.5mm, left=3mm, right=3mm, top=2mm, bottom=2mm, breakable]
The graph shows $p(x) = x^3 - 2x^2 - 5x + 6 = (x+2)(x-1)(x-3)$.
\begin{center}
\begin{tikzpicture}[scale=0.46]
  \lgrid{-3}{4}{-6}{10}
  \begin{scope}
    \clip (-3,-6) rectangle (4,10);
    \draw[very thick, forest, domain=-2.6:3.6, samples=100, smooth]
      plot (\x,{(\x)*(\x)*(\x) - 2*(\x)*(\x) - 5*(\x) + 6});
  \end{scope}
  \fill[charcoal] (-2,0) circle (3pt) node[below left, font=\scriptsize]{$(-2,0)$};
  \fill[charcoal] (1,0) circle (3pt) node[below right, font=\scriptsize]{$(1,0)$};
  \fill[charcoal] (3,0) circle (3pt) node[below right, font=\scriptsize]{$(3,0)$};
  \fill[charcoal] (0,6) circle (3pt) node[right, font=\scriptsize]{$(0,6)$};
\end{tikzpicture}
\end{center}
\begin{enumerate}[leftmargin=1.7em, itemsep=7pt]
  \item The degree is (circle one): \textbf{even} \; / \; \ans{\textbf{odd}}. \quad Its arms: the left arm
        points \ans{down} and the right arm points \ans{up}.
  \item The graph turns \ans{2} times (its turning points).
  \item The \textbf{zeros} are $x=\ans{-2}$, \; $x=\ans{1}$, \; $x=\ans{3}$, and each
        one \textbf{crosses} the axis. The \textbf{$y$-intercept} is $(0,\ \ans{6})$.
  \item Between the zeros the curve makes one hill and one valley. The hill is a
        \textbf{relative \ans{maximum}} and the valley is a \textbf{relative \ans{minimum}}.
\end{enumerate}
\end{tcolorbox}

\begin{tcolorbox}[colback=white, colframe=black!40, title=\textbf{Tier A --- Approaching Proficiency},
    fonttitle=\bfseries, arc=1.5mm, left=3mm, right=3mm, top=2mm, bottom=2mm, breakable]
For \emph{each} graph, give the full feature list, then answer the justification below.
\begin{center}
\begin{tikzpicture}[scale=0.42]
  % p(x) = x^3 - 4x  (odd)
  \begin{scope}
    \lgrid{-3}{3}{-4}{4}
    \begin{scope}
      \clip (-3,-4) rectangle (3,4);
      \draw[very thick, forest, domain=-2.4:2.4, samples=90, smooth]
        plot (\x,{(\x)*(\x)*(\x) - 4*(\x)});
    \end{scope}
    \fill[charcoal] (-2,0) circle (3pt); \fill[charcoal] (0,0) circle (3pt);
    \fill[charcoal] (2,0) circle (3pt);
    \node[charcoal, font=\scriptsize] at (0,-5.2){$p(x)=x^3-4x$};
  \end{scope}
  % r(x) = -x^4 + 4x^2  (even)
  \begin{scope}[xshift=9cm]
    \lgrid{-3}{3}{-3}{5}
    \begin{scope}
      \clip (-3,-3) rectangle (3,5);
      \draw[very thick, forest, domain=-2.3:2.3, samples=100, smooth]
        plot (\x,{-1*(\x)^4 + 4*(\x)*(\x)});
    \end{scope}
    \fill[charcoal] (0,0) circle (3pt); \fill[charcoal] (-2,0) circle (3pt);
    \fill[charcoal] (2,0) circle (3pt);
    \fill[charcoal] (1.414,4) circle (3pt); \fill[charcoal] (-1.414,4) circle (3pt);
    \node[charcoal, font=\scriptsize] at (0,-4.2){$r(x)=-x^4+4x^2$};
  \end{scope}
\end{tikzpicture}
\end{center}
\renewcommand{\arraystretch}{1.5}
\begin{tabularx}{\linewidth}{@{}>{\bfseries}p{0.30\linewidth} X X@{}}
 & \textbf{$p(x)=x^3-4x$} & \textbf{$r(x)=-x^4+4x^2$} \\
\hline
Degree even / odd & \ans{odd} & \ans{even} \\
End behavior (both arms) & \ans{left down, right up} & \ans{both arms down} \\
Number of turning points & \ans{2} & \ans{3} \\
Zeros (cross / touch) & \ans{$-2,0,2$ (all cross)} & \ans{$\pm2$ cross; $0$ touches} \\
Absolute max / min? & \ans{neither} & \ans{abs.\ max $4$; no abs.\ min} \\
Symmetry & \ans{odd} & \ans{even} \\
\end{tabularx}
\vspace{2pt}

\textbf{Justify:} Show algebraically that $p(x)=x^3-4x$ is \emph{odd} by computing $p(-x)$ and
comparing it to $p(x)$. What does that tell you about the graph's symmetry?
\ansline{$p(-x)=(-x)^3-4(-x)=-x^3+4x=-(x^3-4x)=-p(x)$. Since $p(-x)=-p(x)$, $p$ is odd, so its graph has $180^\circ$ rotational symmetry about the origin.}
\end{tcolorbox}

\begin{tcolorbox}[colback=white, colframe=black!40, title=\textbf{Tier E --- Extension},
    fonttitle=\bfseries, arc=1.5mm, left=3mm, right=3mm, top=2mm, bottom=2mm, breakable]
\textbf{Part 1 --- Degree detective.} A polynomial's \emph{left} arm points \emph{up} and its
\emph{right} arm points \emph{down}, and the graph has exactly \textbf{3} turning points.
\begin{enumerate}[label=(\alph*), leftmargin=1.8em, itemsep=4pt]
  \item Is the degree even or odd? \ans{odd} \quad How do the arms tell you? \ans{they point in \emph{opposite} directions, which only happens for odd degree}
  \item Is the leading coefficient positive or negative? \ans{negative} Why? \ans{the right arm points down, so $a<0$}
  \item A degree-$n$ graph has at most $n-1$ turning points. What is the \emph{smallest} degree this
        polynomial could have? \ans{5}
\end{enumerate}
\begin{itemize}[leftmargin=1.4em, nosep]
  \item[] \footnotesize $3$ turning points need degree $\ge 4$; but the degree must be \emph{odd} (opposite arms), so the smallest is $5$.
\end{itemize}

\vspace{4pt}
\textbf{Part 2 --- Open-box model.} A rectangular sheet $10\text{ in}\times 8\text{ in}$ has a square
of side $x$ cut from each corner; the flaps fold up into an open box of volume
\[
  V(x) = x(10-2x)(8-2x).
\]
\begin{center}
\renewcommand{\arraystretch}{1.25}
\begin{tabular}{|c|c|c|c|c|c|c|}
\hline
$x$ & $0$ & $1$ & $1.5$ & $2$ & $3$ & $4$ \\ \hline
$V(x)$ & $0$ & $48$ & $52.5$ & $48$ & $24$ & $0$ \\ \hline
\end{tabular}
\end{center}
\begin{enumerate}[label=(\alph*), leftmargin=1.8em, itemsep=6pt]
  \item The \emph{feasible domain} is $0<x<4$, not all real numbers. Explain why, using the fact that
        $x$ is a length and the flap $8-2x$ must be positive.
        \ansline{$x$ is a cut length, so $x>0$; and $8-2x>0$ forces $x<4$. Outside $0<x<4$ a side would be zero or negative, so no real box exists.}
  \item From the table, which $x$ gives the greatest volume, and what is that greatest volume?
        \ans{$x=1.5$ in gives the greatest volume, $V=52.5$ in$^3$.}
  \item Fully expanded, $V$ has zeros at $x=0,\ 4,\ 5$. Which of these are \emph{meaningful} for the
        box, and what does each meaningful zero represent?
        \ansline{$x=0$ and $x=4$ are meaningful --- they bound the feasible domain and each gives $V=0$ (no box): $x=0$ means no corner is cut, $x=4$ cuts the $8$-in side away entirely. $x=5$ is outside $0<x<4$, so it is not meaningful.}
\end{enumerate}
\end{tcolorbox}

\begin{teachernote}
Tier A's key idea is that ``relative'' and ``absolute'' can differ: the odd-degree $p$ has two
turning points but \emph{no} absolute extrema (arms run to $\pm\infty$), while the even-degree $r$'s
two peaks \emph{are} its absolute maximum. In Tier E Part 1(c), the common error is answering ``4''
--- remind students the degree must stay odd, so it jumps to 5. Accept $x\approx1.5$ in Part 2(b);
the true maximizing cut is near $1.47$ in.
\end{teachernote}

\end{document}
